\section{Evolution of Industrial Object Detection and Counting}
The automation of object detection and counting on industrial conveyor belts has evolved significantly from traditional computer vision to deep learning paradigms. Early approaches relying on handcrafted features such as Histogram of Oriented Gradients (HOG) or Scale-Invariant Feature Transform (SIFT) combined with Support Vector Machines (SVM) proved highly susceptible to variable lighting, motion blur, and complex backgrounds. The shift toward Convolutional Neural Networks (CNNs), particularly the You Only Look Once (YOLO) architecture family, has established a new baseline for industrial applications.\\

A foundational study by \cite{jiang2021accurate} demonstrated the viability of a YOLOv4-based system for fastener counting, utilizing K-Means anchor redesign and transfer learning to achieve a remarkable 96.43\% mAP@0.5. However, this system was constrained to an inference speed of only 11 FPS on a high-end GTX 1080 Ti GPU, rendering it unsuitable for real-time edge deployment. Similarly, comprehensive evaluations on industrial datasets, such as the UDD (Industrial Recycling) benchmark by \cite{Messai_2026}, confirm that while larger models like YOLOv8x achieve superior performance (78.3\% mAP@0.5) for small, dense objects, they still suffer from performance degradation when objects are heavily occluded or exhibit extreme scale variability. These findings underscore a persistent challenge in industrial vision: balancing high-accuracy small object detection with the stringent latency requirements of real-time conveyor systems.

\section{Attention Mechanisms for Small and Occluded Objects}
To mitigate the loss of feature representation in small and overlapping objects, attention mechanisms have been widely integrated into object detection pipelines. However, many traditional attention modules introduce prohibitive computational overhead or suffer from dimensionality reduction side effects. For instance, \cite{Liu2025YOLO} integrated a BotNet attention mechanism into YOLOv5s for conveyor belt tear detection, replacing standard convolutions with Multi-Head Self-Attention (MHSA). While this improved mAP by 2.9\%, it inherently added Transformer-based computational overhead and did not adequately address severe occlusion.\\

In contrast, the Efficient Multi-Scale Attention (EMA) module proposed by \cite{Ouyang_2023} offers a more optimal trade-off for industrial edge deployment. EMA utilizes multi-scale parallel subnetworks (1$\times$1 and 3$\times$3 branches) and cross-spatial learning via matrix dot-product fusion. Crucially, it operates \textit{without} dimensionality reduction, preserving full channel information and capturing fine-grained, pixel-level relationships necessary for distinguishing overlapping fasteners. While alternative lightweight attention mechanisms like Fast Window Attention (FWA) with DReLu activation \cite{li2025lightweightbackbonenetworksrequire} offer adaptive key caching and faster inference, they tend to lose fine-grained details in high-resolution industrial images. Therefore, EMA presents a superior foundational mechanism for integration into modern YOLO architectures, directly inspiring the novel \textbf{EMAc2f} module proposed in this work.

\section{Lightweight Architectures for Edge Deployment}
Deploying robust computer vision models on resource-constrained edge devices, such as the NVIDIA Jetson series, necessitates architectural optimization. Standard CNN backbones contain significant feature map redundancy. \cite{han2020ghostnetfeaturescheapoperations} addressed this by introducing GhostNet, which generates redundant feature maps through cheap linear operations rather than expensive standard convolutions. This approach yields a theoretical $\sim$2$\times$ speedup and parameter compression with minimal accuracy degradation, making it an ideal backbone candidate for edge-based industrial vision.\\

The practicalities of edge deployment in industrial settings have also been explored through multi-modal approaches. \cite{lim2025realtimemultimodalobjectdetection} developed a real-time system combining Time-of-Flight (ToF) and RGB cameras on a Jetson AGX Orin, utilizing unsupervised density-based detection (HDBSCAN and Gaussian KDE). While their system proved robust to low-light conditions and closely clumped objects, the KDE pre-processing was computationally intensive and prone to disruption by extraneous objects. By contrast, integrating a lightweight deep learning model (GhostNet + YOLO11) with RGB-D sensors (e.g., Intel RealSense) offers a more scalable, accurate, and end-to-end trainable alternative for continuous fastener counting.

\section{Temporal Consistency and Tracking in Industrial Vision}
A critical, often overlooked limitation in industrial counting systems is the reliance on frame-wise detection, which produces fluctuating predictions due to motion blur, transient occlusions, and illumination changes. \cite{park2026twostagedetectiontrackingframeworkstable} addressed this in agricultural quality inspection by implementing a two-stage framework: YOLOv8 for detection followed by ByteTrack for multi-object tracking, coupled with track-level majority voting for classification. This methodology significantly improved temporal stability and reduced prediction oscillation. Although initially applied to a different domain, this tracking paradigm provides a strong architectural inspiration for the proposed system, where persistent tracking IDs can be leveraged to prevent the double-counting of fasteners as they traverse the conveyor belt.

\section{Identified Research Gaps and Proposed Contributions}
Synthesizing the reviewed literature reveals several critical gaps in current industrial fastener counting systems:
\begin{enumerate}
    \item \textbf{Speed-Accuracy Trade-off:} Existing high-accuracy models (e.g., YOLOv4, YOLOv8x) are too computationally heavy for real-time edge deployment, while lightweight models sacrifice critical small-object detection capabilities.
    \item \textbf{Occlusion and Scale Handling:} Overlapping fasteners and extreme scale variability frequently cause missed detections in standard YOLO architectures due to inadequate multi-scale feature fusion.
    \item \textbf{Industrial Domain Gap:} Models trained on general-purpose datasets often fail to generalize to the specific visual characteristics, lighting, and motion dynamics of industrial conveyor belts.
\end{enumerate}

\textbf{Proposed Contributions:} To bridge these gaps, this project proposes \textbf{YOLO11-GED}, a novel, full-stack architecture tailored for edge deployment. The unique contributions of this work, directly responding to the identified literature gaps, include:
\begin{itemize}
    \item \textbf{Novel EMAc2f Module:} A custom integration of the EMA attention mechanism into the YOLO C2f structure, enabling multi-scale attention without dimensionality reduction to robustly resolve overlapping fasteners.
    \item \textbf{GhostNet Backbone Integration:} Replacement of the standard heavy backbone with GhostNet to achieve significant computational speedup and parameter reduction, specifically optimized for NVIDIA Jetson platforms.
    \item \textbf{DySample Dynamic Upsampling:} Utilization of dynamic, content-aware upsampling to enhance small feature map recovery, directly addressing the scale variability of industrial parts.
    \item \textbf{Full-Stack Edge Deployment Pipeline:} A complete, practical implementation integrating the optimized YOLO11-GED model with an Intel RealSense camera on an NVIDIA Jetson platform, bridging the gap between theoretical model design and real-world industrial application.
\end{itemize}